The model was developed in an existing Agent-Based Modelling Software Python library called \textit{MESA} \citep{Mesa2025}, a methodological choice strongly supported by recent reviews of computational tools in solid waste management \citep{TianReview2024, Ma2023}. Its architecture consisted of one single ABM environment that contained 7 distinct \texttt{BarangayAgent} objects, which in turn contained their respective populations of \texttt{HouseholdAgent} objects.

The simulation utilized a single, discrete ABM environment to contain its operational entities and manage its temporal structure. Time progression was segmented into recurring quarterly loops, establishing the fundamental unit of analysis and adaptation for the governance agents. The architecture was strictly hierarchical: it contained seven distinct \texttt{BarangayAgent} objects in the upper layer, which encapsulated their respective populations of \texttt{HouseholdAgent} objects in the lower layer. This multi-level approach was essential for accurately capturing the flow of mandates and the heterogeneity between different governance units in socio-environmental systems \citep{Brugiere2022}.

\begin{figure}[h]
    \centering
    % Ensure you have uploaded the file 'multi-agent.png' to your project
    \includegraphics[width=0.8\linewidth]{multi-agent.png}
    \caption{Multi-Level Agent-Based Modeling Diagram}
    \label{fig:multi_level_abm}
\end{figure}

% Define a custom bullet for tables to save space and look cleaner
\newcommand{\tabitem}{~~\llap{\textbullet}~~}

\begin{table}[htbp]
    \centering
    \small
    \caption{Agent Role Mapping and Architecture}
    \label{tab:agent_roles_v2}
    \renewcommand{\arraystretch}{1.5} 
    \begin{tabularx}{\textwidth}{ 
        >{\RaggedRight\hsize=0.15\hsize}X 
        >{\RaggedRight\hsize=0.12\hsize}X 
        >{\RaggedRight\hsize=0.29\hsize}X 
        >{\RaggedRight\hsize=0.22\hsize}X 
        >{\RaggedRight\hsize=0.22\hsize}X 
    }
        \toprule
        \textbf{Agent} & \textbf{Type} & \textbf{Input / Observation} & \textbf{Action Space} & \textbf{Objective} \\
        \midrule
        \textbf{Household} & ABM (Env.) & 
        \tabitem Policy: Incentive \& Fine amounts \newline
        \tabitem Social: Neighbor compliance \newline & 
        Decide: Segregate or Do Not Segregate (Binary) & 
        Maximize Utility: Minimize ``Cost of Effort'' vs. Fear of Penalty. \\
        \textbf{Barangay} & ABM (Impl.) & 
        \tabitem Budget: Quarterly funds \newline
        \tabitem Population: Count \& income profile & 
        Execute: Deploy Enforcers or distribute rewards & 
        Report State: Aggregate data and transmit to Mayor Agent. \\
        \textbf{Mayor (Municipal)} & DRL / AI (Brain) & 
        \tabitem State: Compliance of 7 Barangays \newline
        \tabitem Fiscal: Remaining Budget \newline
        \tabitem Political: Political Capital & 
        Allocate: Distribute budget to (1) IEC, (2) Incentives, (3) Enforcement. & 
        Maximize Reward: Balance compliance with sustainability. \\
        \textbf{Enforcement} & ABM (Ops.) & 
        \tabitem Target: Assigned Barangay \newline
        \tabitem Constraint: Working hours & 
        Inspect: Randomly check households & 
        Detect: Trigger ``Perceived Probability of Detection''. \\
        \bottomrule
    \end{tabularx}
\end{table}

The \textit{BarangayAgents} embodied the frontline implementation layer, responsible for baseline fiscal management, policy execution (e.g., establishing waste segregation mandates), and personnel management (deployment of enforcement staff) using their own separate funds. This aligned with the decentralized governance structure where local units were the primary implementers of environmental law \citep{Nishimura2022}. The \textit{BarangayAgents} were the key aggregators for the feedback mechanism, utilizing the combined compliance rate as an observational State and prescriptive Reward signal transmitted to the central \textbf{Mayor Agent} (the Deep Reinforcement Learning mechanism) \citep{Jimenez2025}.

The \textit{EnforcementAgents} functioned as the operational execution arm, tasked with the stochastic monitoring of compliance. Constrained by the budget allocation, these agents conducted random inspections to verify waste segregation at the source, effectively operationalizing the "No Segregation, No Collection" policy \citep{Badua2022}. Their presence directly manipulated the "perceived intensity of enforcement" variable within the simulation, transforming abstract policy mandates into tangible risks for non-compliant households \citep{Carpio2025}. This agent layer introduced the critical dynamic of detection probability: while higher enforcement density significantly deterred non-compliance ("green criminology"), it incurred higher operational costs, forcing the system to balance fiscal sustainability with strict environmental policing \citep{Chen2023, Ibanez2022}.

The \textit{HouseholdAgents} represented the citizenry and were geographically and administratively segregated by their assigned barangay. Household behavior was modeled through a process of rational bounded observation, where agents assessed three primary exogenous factors: the clarity and strictness of local policies, the perceived intensity of enforcement (driven by the \textit{EnforcementAgents}), and the compliance behavior of neighboring households (social norms) \citep{Meng2018, Liao2024}. This assessment culminated in an independent segregation decision, which dictated the agent's contribution to the collective compliance level. The quarterly feedback loop aggregated these decisions into seven discrete \textit{BarangayComplianceRate} values, which were then transmitted back to the Mayor Agent, closing the adaptive system loop and facilitating the study of municipal policy response to citizen behavior \citep{Ceschi2021, Fontaine2024}.

\subsection{The Household Agent: Behavioral and Computational Architecture}

The \texttt{HouseholdAgent} is the foundational computational node of the simulation. Rather than acting as a static data point, each household acts as a continuous finite state machine. To ensure the reader can follow the complex mechanics of this agent, this section is divided into three distinct phases: (1) the theoretical mathematical framework governing behavior, (2) the computational data structures and spatial algorithms that execute this theory, and (3) the validation of these parameters via Global Sensitivity Analysis.

\subsubsection{Theoretical Foundation: The TPB Framework and Utility Architecture}

The core decision logic of the \texttt{HouseholdAgent} is governed by a dynamic utility function grounded in the Theory of Planned Behavior (TPB). This theory posits that an individual's intention to perform a behavior is a function of their Attitude ($A$), Subjective Norms ($SN$), and Perceived Behavioral Control ($PBC$) \citep{Ceschi2021}. In contrast to static behavioral models, this framework incorporates time-variant weights ($w(t)$) for psychological constructs, allowing agent behavior to evolve non-linearly \citep{Taraghi2025, Ma2023}. The baseline utility of segregation ($U_{\mathrm{segregate}}$) is defined as:

\begin{equation}
    U_{\mathrm{segregate}} = (w_A(t) A) + (w_{SN}(t) SN_{\mathrm{local}}) + (w_{PBC}(t)  PBC) - C_{\mathrm{Net}} + \epsilon
    \label{eq:utility_function}
\end{equation}
\addequation{Household Utility of Segregation Function}

Where the decision to segregate is a stochastic process driven by the following parameters:

\begin{description}
    \item[$w_A(t) A$] \textbf{Attitude Component:} Represents the internal motivation. $A$ is the dynamic sentiment score (boosted by IEC campaigns), while $w_A$ is the psychological weight the agent assigns to personal belief.
    \item[$w_{SN}(t) SN_{\mathrm{local}}$] \textbf{Social Norm Component:} The influence of the community. $SN_{\mathrm{local}}$ is the local compliance rate of neighbors, and $w_{SN}$ is the susceptibility to peer pressure, which scales non-linearly via the State-Dependent Piecewise Function.
    \item[$w_{PBC}(t) PBC$] \textbf{Perceived Behavioral Control:} The agent's belief in their capability to segregate, representing access to resources (e.g., bins or collection points).
    \item[$C_{\mathrm{Net}}$] \textbf{Net Cost of Effort:} The physical and economic friction of segregation ($c_{\text{effort}}$) minus any active financial incentives ($I$) provided by the Mayor Agent ($C_{\mathrm{Net}} = c_{\text{effort}} - I$).
    \item[$\epsilon$] \textbf{Stochastic Noise:} A Gaussian noise term $(\mu=0.0, \sigma=0.05)$ modeling the inherent unpredictability of human decision-making.
\end{description}

In the utility function, a household's internal disposition toward waste segregation is represented by two distinct variables: the Attitude score ($A$) and its corresponding psychological weight ($w_A$). The weight ($w_A$) represents the intrinsic importance an individual places on their own personal beliefs relative to external factors like peer pressure or penalties. This weight remains a relatively stable baseline trait for each agent.

Conversely, the Attitude score itself ($A$) is highly dynamic and represents the agent's current level of positive sentiment toward segregation. The value of $A$ increases directly in response to municipal Information, Education, and Communication (IEC) campaigns allocated by the Mayor Agent. Following an IEC boost, the Attitude score ($A$) decays stochastically over time, modeling the psychological reality that the impact of public awareness campaigns fades without continuous reinforcement..

A critical computational innovation in this model is the derivation of $C_{\mathrm{Net}}$, the net behavioral cost, which translates abstract psychological friction into a computable scalar:

\begin{equation}
    C_{\mathrm{Net}} = C_{\mathrm{Effort}} + (\gamma C_{\mathrm{Monetary}}) - (\gamma I) - (\gamma F P_{\mathrm{Detection}})
    \label{eq:net_disutility}
\end{equation}
\addequation{Net Disutility of Segregation Function}

This cost-benefit sub-model integrates both physical and financial barriers:
\begin{itemize}
    \item \textbf{$C_{\mathrm{Effort}}$}: The physical hassle associated with washing, sorting, and storage of waste.
    \item \textbf{$C_{\mathrm{Monetary}}$}: Tangible financial expenses, such as procuring color-coded bins.
    \item \textbf{$\gamma$ (Income Sensitivity)}: A weighting factor derived from household income level; for lower-income households, $\gamma > 1$ as financial levers carry greater psychological weight, whereas $\gamma < 1$ for higher-income households.
    \item \textbf{$I$, $F$, and $P_{\mathrm{Detection}}$}: The objective magnitudes of monetary incentives, punitive fines, and detection probability. \textit{(Note: These external variables are dynamically controlled by macro-level municipal interventions, which will be extensively detailed in Section 3.3.4: The Mayor Agent).}
\end{itemize}

The model also captures complex psychological feedbacks such as ``psychological reactance.'' While municipal investments generally improve Attitude ($A$), if enforcement intensity crosses a specific threshold, Attitude paradoxically decreases, reflecting resistance to coercive mandates \citep{LoanBalanay2023}.

\subsubsection{Modeling Behavioral Inertia: Social Shields and Norm Internalization}

Standard Agent-Based Models often suffer from hyper-reactivity, exhibiting unrealistic behavioral decay where compliance collapses immediately upon the cessation of external stimuli. To accurately simulate the real-world stability of established habits, this study mathematically models ``Cultural Inertia''---where a behavior transitions from a calculated, incentive-driven decision to an internalized social habit \citep{Corcoran2020}.

To achieve this without relying on complex path-dependent variables, the model implements a ``Social Norm Shield'' using a state-dependent piecewise function. This shield categorizes community behavior into distinct cultural states based on current localized compliance thresholds, effectively shifting the behavioral driver from transient attitudes to structural peer pressure \citep{Andre2021}. The model operates through the following programmatic states:

\begin{itemize}
    \item \textbf{Established Culture (The 70\% Tipping Point):} Once a localized node breaches the 70\% compliance threshold, the behavior transitions from incentive-driven to norm-driven. The simulation executes a deterministic hard-lock on the Subjective Norm weight ($w_{sn} = 0.90$). This state represents strong cultural inertia and reflects ``critical mass'' theory, where minority behaviors cascade into majority conventions that actively resist degradation \citep{Centola2018, Nyborg2024}.
    \item \textbf{Transitional State (The 40\% Bayanihan Floor):} When compliance is between 40\% and 69\%, the system models an emergent community spirit. To prevent total mathematical collapse during the DRL agent's early exploration phase, an absolute floor is implemented. Even if enforcement ceases entirely, spatial peer pressure prevents the social norm weight from dropping below a baseline. This simulates the internalized ``moral obligation'' and residual cultural memory identified in regional waste management studies \citep{Nishimura2022}.
    \item \textbf{Systemic Failure (Below 40\%):} If compliance fails to reach or drops below 40\%, the social shield collapses. At this stage, structural peer pressure is virtually non-existent, and the community relies entirely on external enforcement and transient attitudes to drive compliance.
\end{itemize}

Mathematically, this mechanism replaces standard dynamic decay with a state-dependent piecewise function that locks or floors the Subjective Norm variable based on the current environment:

\begin{equation}
    w_{sn}(t) = 
    \begin{cases} 
        0.90 & \text{if } \text{Compliance}_{\mathrm{local}} \ge 0.70 \text{ (Established Culture)} \\
        \max(w_{sn}(t_0), 0.40) & \text{if } 0.40 \le \text{Compliance}_{\mathrm{local}} < 0.70 \text{ (Transitional State)} \\
        w_{sn}(t_0) & \text{otherwise (Systemic Failure)}
    \end{cases}
\end{equation}
\addequation{State-Dependent Piecewise Function for Subjective Norms}

\subsubsection{Computational Implementation and State Transitions}

While the preceding equations establish the theoretical bounds of the agent, the \texttt{HouseholdAgent} requires a computational architecture to process these variables dynamically. In the ABM framework, the agent executes a continuous evaluation loop (the \texttt{step()} function) at every discrete time step ($t$). 

\subparagraph{Part 1: State Vector Initialization:} 
Upon instantiation, each agent is assigned a fixed demographic parameter (Income Level) and a continuously updating psychological state vector $S_{hh} = \langle A, SN, PBC \rangle$. To prevent synthetic uniformity, the initial state vector is populated using stochastic Gaussian sampling bounded by the agent's assigned behavioral profile \citep{Fischer2025}. For instance, agents initialized as "Compliant" derive their base attitude from $A \sim \mathcal{U}(0.75, 0.90)$, whereas "Non-Compliant" agents are seeded from $A \sim \mathcal{U}(0.10, 0.40)$.

\subparagraph{Part 2: Spatial Queries (Moore Neighborhood Algorithm):} 
The calculation of Social Norms ($SN_{\mathrm{local}}$) relies on spatial data. During execution, the agent utilizes a Moore Neighborhood search algorithm with a radius of $r=1$, querying adjacent grid cells in $O(k)$ time, where $k \le 8$ cells. The local compliance ratio $R_{local}$ is computed as:
\begin{equation}
    R_{local} = \frac{\sum_{i=1}^{k} \mathbb{I}(\text{neighbor}_i.\text{is\_compliant})}{k}
\end{equation}
If $R_{local}$ crosses the predefined social tipping point of $0.25$, the agent's internal $SN$ weight undergoes an asymmetric scalar multiplication ($SN_{new} = SN_{old} \times 1.30$), triggering a rapid, non-linear computational cascade in community compliance \citep{Centola2018}.

\subparagraph{Part 3: Event-Driven Overrides (\texttt{get\_fined}):} 
The \texttt{HouseholdAgent} exposes public methods for event-driven interrupts triggered by spatial collisions with an \texttt{EnforcementAgent}. When \texttt{get\_fined(amount)} is invoked, an economic penalty is applied to the agent's utility, scaled dynamically by the $\gamma$ multiplier. Furthermore, if the fine exceeds \textpeso 500, a probabilistic "Fear Factor" algorithm executes: a 60\% probability that the agent's internal state is forcefully overridden to \texttt{is\_compliant = True} for a temporal window. However, this applies a permanent flat deduction to the Attitude variable ($\Delta A = -0.10$), creating a computational tradeoff between short-term forced compliance and long-term intrinsic decay \citep{Hertweck2023}.

\paragraph{Part 4: Algorithmic Execution of Household Decisions}

The integration of the mathematical utility functions, the protective social shields, and the stochastic noise filters is executed daily by each household. To synthesize these complex operations into an interpretable sequence, Algorithm \ref{alg:household_decision_english} details the exact step-by-step heuristic process the agent undergoes during every simulation tick.

\begin{algorithm}[H]
\caption{The Household's Daily Decision Process (To Segregate or Not?)}
\label{alg:household_decision_english}
\begin{algorithmic}[1]
\State \textbf{Input:} Household Income, Baseline Motivation, Peer Pressure sensitivity, and Base Effort.
\Statex
\State \textbf{Step 1: Evaluate Peer Pressure \& The Habit Shield}
\State \parbox[t]{0.9\linewidth}{$\rightarrow$ Calculate the local segregation rate by checking immediate neighbors.\strut}
\If{the neighborhood is doing worse than the household expected}
    \State \parbox[t]{0.88\linewidth}{$\rightarrow$ \textit{Apply Habit Shield:} If the neighborhood is highly compliant ($>70\%$), heavily resist the negative pressure. If average ($>40\%$), resist slightly. If failing, instantly adopt the negative behavior.\strut}
\Else
    \State \parbox[t]{0.88\linewidth}{$\rightarrow$ Immediately absorb positive social pressure to fit in.\strut}
\EndIf
\Statex
\State \textbf{Step 2: Update Motivation \& Calculate Net Effort}
\State \parbox[t]{0.9\linewidth}{$\rightarrow$ Apply natural motivation decay, buffered by strong community habits, and apply motivation boosts from recent Educational Campaigns.\strut}
\State \parbox[t]{0.9\linewidth}{$\rightarrow$ Calculate the "Net Effort" by weighing the physical hassle against the financial impact of municipal fines and rewards (scaled by the household's income level).\strut}
\State \parbox[t]{0.9\linewidth}{$\rightarrow$ Add a fraction of random human unpredictability to the thought process.\strut}
\Statex
\State \textbf{Step 3: Execute the Final Choice}
\If{the overall benefits and social pressures outweigh the effort}
    \State \parbox[t]{0.88\linewidth}{$\rightarrow$ \textit{Stochastic Human Error Check:} Roll a 1\% chance of accidental failure. If passed, successfully segregate waste.\strut}
\Else
    \State \parbox[t]{0.88\linewidth}{$\rightarrow$ The effort is too high or motivation is too low: Do not segregate.\strut}
\EndIf
\end{algorithmic}
\end{algorithm}


\subsubsection{Validation of Household Mechanics: Global Sensitivity Analysis}

Agent-Based Models (ABMs) are highly effective at simulating complex systems, but they frequently suffer from parameter uncertainty, particularly when dealing with unobservable psychological constructs \citep{McCulloch2022}. To rigorously validate the structural integrity of the Household Agent's programming and quantify this uncertainty, a variance-based Global Sensitivity Analysis (GSA) was conducted using the Sobol method \citep{Sobol2001}. 

Unlike local sensitivity analyses that alter one variable at a time, the Sobol method explores the entire multidimensional parameter space simultaneously. This approach decomposes the variance of the model's primary macro-level outputs—specifically, the final Global Compliance rate—into fractions directly attributable to the individual Theory of Planned Behavior (TPB) parameters and their complex, non-linear interactions \citep{ZhengAI2022}. 

\begin{algorithm}[H]
\caption{Global Sensitivity Analysis (Stress-Testing the Simulation)}
\label{alg:sobol_gsa_english}
\begin{algorithmic}[1]
\State \textbf{Starting Information Needed:} 
\State \parbox[t]{0.9\linewidth}{$\rightarrow$ The boundaries for the psychological traits and a target of 384 parallel simulation runs.\strut}
\Statex
\State \textbf{Step 1: Define the "What-If" Variables}
\State \parbox[t]{0.9\linewidth}{$\rightarrow$ Set the minimum and maximum limits for the 5 core behavioral drivers: Personal Motivation ($w_a$), Peer Pressure ($w_{sn}$), Rule Difficulty ($w_{pbc}$), Physical Effort ($C_{\mathrm{Effort}}$), and Motivation Decay.\strut}
\Statex
\State \textbf{Step 2: Generate the Test Matrix (Saltelli Method)}
\State \parbox[t]{0.9\linewidth}{$\rightarrow$ Mathematically scramble and mix these 5 traits together to create a massive list of unique citizen profiles (e.g., highly conformist communities vs. highly independent communities).\strut}
\Statex
\State \textbf{Step 3: Run the Simulations (Headless Mode)}
\For{every unique citizen profile in the test matrix}
    \State \parbox[t]{0.88\linewidth}{$\rightarrow$ Inject these specific psychological traits into the simulation's Household Agents.\strut}
    \State \parbox[t]{0.88\linewidth}{$\rightarrow$ Fast-forward the simulation to the end of the 3-year term using the standard "Status Quo" policy.\strut}
    \State \parbox[t]{0.88\linewidth}{$\rightarrow$ Record the municipality's final Waste Segregation Rate and the Mayor's Political Capital score.\strut}
\EndFor
\Statex
\State \textbf{Step 4: Calculate the Impact Scores (Sobol Indices)}
\State \parbox[t]{0.9\linewidth}{$\rightarrow$ Compute the \textbf{First-Order Index ($S_1$)}: This reveals the "Direct Impact" (which single trait changes the final segregation rate on its own).\strut}
\State \parbox[t]{0.9\linewidth}{$\rightarrow$ Compute the \textbf{Total-Order Index ($S_T$)}: This reveals the "Ripple Effect" (which trait changes outcomes when it interacts and combines with other psychological traits).\strut}
\end{algorithmic}
\end{algorithm}

The parameter space was deliberately bounded to reflect plausible ranges of human behavioral drivers derived from existing socio-psychological literature \citep{Meng2018}. The psychological weights ($w_a, w_{sn}, w_{pbc}$) were uniformly bounded between $[0.1, 0.6]$. This wide yet constrained range ensures that no single psychological factor deterministically overpowers the algorithmic decision-making without empirical justification. Additionally, operational constraints such as Base Effort and Attitude Decay were bounded at $[0.1, 0.5]$ and $[0.001, 0.015]$, respectively. 

By calculating the $S_1$ and $S_T$ indices, the analysis isolates which specific computational parameters most heavily dictate systemic outcomes. Identifying these highly sensitive parameters mathematically validates whether the emergent behavioral dynamics of the ABM align with established psychological theories of waste segregation, ensuring the subsequent Deep Reinforcement Learning (DRL) agent is optimizing against a structurally sound and ecologically valid virtual population \citep{TianReview2024}.

\subsection{The Enforcement Agent: Spatial Dynamics and Punitive Actuation}

Operating under the directive of both the local \texttt{BarangayAgent} and the municipal Mayor, the \texttt{EnforcementAgent} serves as the punitive actuator of the simulation. Computationally, this agent is the physical manifestation of the $F$ (Fine) and $P_{\mathrm{Detection}}$ variables found within the Household's utility function.

\paragraph{Part 1: Jurisdictional Typology and Temporal Constraints}

To accurately reflect the decentralized governance of Philippine Solid Waste Management \citep{Nishimura2022}, \texttt{EnforcementAgent} objects are dynamically instantiated under two distinct operational profiles, each with unique computational constraints:

\begin{enumerate}
    \item \textbf{Micro-Enforcement (Barangay Tanods):} Financed by the baseline barangay budget, these persistent agents simulate local village watchmen. Because real-world Tanods are burdened with multiple community duties beyond waste management, their \texttt{step()} function executes a stochastic distraction filter. At the start of each daily tick, a pseudorandom float determines their availability; computationally, there is a 90\% probability (`rand() > 0.10`) that the agent yields its execution thread to other assumed duties and skips garbage monitoring entirely. Their daily apprehension capacity is strictly bottlenecked to a maximum of one non-compliant household.
    
    \item \textbf{Macro-Enforcement (LGU Strike Force):} When the Deep Reinforcement Learning (DRL) Mayor allocates municipal funds to enforcement, it spawns dedicated LGU Inspectors. These agents bypass the stochastic distraction filter, executing targeted patrols every day. They possess a high daily apprehension capacity (up to 25 households) and issue severe municipal citations ($F = \text{\textpeso}1,000$). To prevent memory leaks and simulate the logistical realities of temporary task forces, these objects are bound by a strict 30-day temporal contract. The agent decrements a \texttt{contract\_days} integer daily; upon reaching zero, the object executes a self-destruct sequence (\texttt{model.grid.remove\_agent(self)}), forcing the AI Mayor to continuously justify their budgetary upkeep \citep{TianReview2024}.
\end{enumerate}

\paragraph{Part 2: Spatial Computing and Bounded Rationality}

The physical environment of the Agent-Based Model (ABM) is constructed as a two-dimensional discrete spatial grid. Unlike centralized, equation-based mathematical models that assume uniform mixing or omniscient actors, this spatial formulation enforces strict topographical limitations \citep{Meng2018}. 

Computationally, Enforcement Agents do not possess an $O(1)$ global lookup table of all non-compliant citizens. Instead, they operate under bounded rationality. They must execute localized spatial queries within a defined line-of-sight (e.g., a Moore neighborhood of $Radius = 10$). This architectural design accurately simulates the limited monitoring capacity of real-world enforcers who can only penalize violations physically observed within their immediate vicinity \citep{Dalugdog2021}.

The algorithmic movement of the enforcer relies on a finite state machine bifurcated into two behavioral states:
\begin{itemize}
    \item \textbf{State $S_0$ (Patrol Mode):} The agent conducts a stochastic spatial random-walk, moving to adjacent coordinates. This represents the routine, unguided neighborhood rounds conducted when no immediate violations are visible.
    \item \textbf{State $S_1$ (Chase Mode):} If the $O(r^2)$ spatial query detects a non-compliant \texttt{HouseholdAgent} within the scanning radius, the enforcer dynamically calculates a directed spatial vector toward the closest violator to execute an intervention \citep{Zhao2022}.
\end{itemize}

\paragraph{Part 3: Algorithmic Execution of Punitive Patrols}

During each simulation step, the deployed agent employs this randomized monitoring strategy, selecting a discrete sample of households up to its assigned capacity limit. If a household is flagged as non-compliant, the enforcer invokes the household's \texttt{get\_fined()} method, triggering the psychological and economic mutations detailed in Section 3.3.1. 

Algorithm \ref{alg:enforcement_patrol_english} details the step-by-step heuristic process the enforcement agent undergoes during its daily patrol sequence.

\vspace{0.5cm} 
\begin{algorithm}[H]
\caption{The Enforcer's Daily Computational Patrol Process}
\label{alg:enforcement_patrol_english}
\begin{algorithmic}[1]
\State \textbf{Starting Information Needed:} 
\State \parbox[t]{0.9\linewidth}{$\rightarrow$ The Enforcer's current grid coordinate, a spatial ``Field of Vision'' (radius = 10 blocks), and their daily quota capacity.\strut}
\Statex
\State \textbf{Step 1: Memory Management (Preventing Infinite Loops)}
\State \parbox[t]{0.9\linewidth}{$\rightarrow$ Add the immediately surrounding households to a ``Visited'' memory set. This prevents the computational logic from freezing or getting stuck inspecting the exact same houses on consecutive ticks.\strut}
\Statex
\State \textbf{Step 2: Spatial Query and Movement Logic}
\State \parbox[t]{0.9\linewidth}{$\rightarrow$ Scan the entire 10-block Field of Vision for unvisited \texttt{HouseholdAgents}.\strut}
\State \parbox[t]{0.9\linewidth}{$\rightarrow$ Filter the queried households to isolate those whose current state is ``Non-Compliant.''\strut}
\Statex
\If{there is at least one Non-Compliant household in sight}
    \State \parbox[t]{0.88\linewidth}{$\rightarrow$ \textit{CHASE MODE:} Calculate the shortest path to the violator closest to the Enforcer's current coordinate.\strut}
    \State \parbox[t]{0.88\linewidth}{$\rightarrow$ Move the Enforcer one step along that directed spatial trajectory.\strut}
\Else
    \State \parbox[t]{0.88\linewidth}{$\rightarrow$ \textit{PATROL MODE:} If no active violations are visible, generate a random legal coordinate and take one step in that direction to continue exploring the grid.\strut}
\EndIf
\Statex
\State \textbf{Step 3: Enforcement Execution (The Punishment Protocol)}
\State \parbox[t]{0.9\linewidth}{$\rightarrow$ Establish a tight ``Catch Zone'' around the Enforcer's new location.\strut}
\For{every household physically located inside the Catch Zone}
    \If{the household is actively Non-Compliant}
        \State \parbox[t]{0.85\linewidth}{$\rightarrow$ Trigger the household's penalty function (decreasing their utility, lowering their attitude, and forcing temporary compliance).\strut}
        \State \parbox[t]{0.85\linewidth}{$\rightarrow$ Log the financial penalty into the municipality's overall collected fines ledger.\strut}
    \EndIf
\EndFor
\end{algorithmic}
\end{algorithm}

\paragraph{Part 4: Systemic Feedback and the Probability of Detection}

Crucially, the physical presence of these agents generates a systemic macro-level feedback loop. While the enforcers act strictly locally, their aggregate presence dynamically alters the environment. Following the daily deployment phase, the simulation calculates the barangay's macro-level Perceived Probability of Detection ($P_{\mathrm{Detection}}$) by taking the ratio of the total inspection capacity ($C_{\mathrm{Quarterly}}$) of all currently active enforcers against the total household population:

\begin{equation}
    P_{\mathrm{Detection}} = \min\left(1.0, \frac{\sum C_{\mathrm{Quarterly}}}{N_{\mathrm{Households}}}\right)
    \label{eq:prob_detection}
\end{equation}
\addequation{Macro-level Perceived Probability of Detection}

This updated scalar is broadcasted back to the \texttt{HouseholdAgent} objects, dynamically increasing the anticipated financial risk within their utility calculations. 

This spatial limitation profoundly reflects the logistical realities of municipal governance. Local Government Units (LGUs) rarely have the manpower to achieve $P_{\mathrm{Detection}} = 1.0$ \citep{Abordo2025}. Consequently, the probability of a household being caught is not absolute; it is a highly localized phenomenon dependent on spatial proximity and the enforcer's stochastic patrol route. This mechanism organically generates the dynamic, fluctuating risk profile experienced by citizens, ensuring that the deterrent effect of punitive policies is realistically constrained by the physical capacity of the simulated law enforcement \citep{Jimenez2025}.

\subsection{The Barangay Agent: Meso-Level Aggregation and Fiscal Execution}

The \texttt{BarangayAgent} serves as the critical meso-level administrative node within the computational architecture. It bridges the gap between the micro-level behavioral shifts of the \texttt{HouseholdAgent} and the macro-level policy interventions of the \texttt{MayorAgent}. Rather than acting as a spatial entity that moves across the grid, the Barangay Agent functions as a localized environment manager for a specific subset of households. 

To systematically detail its mechanics, this section is divided into three parts: (1) State initialization and data aggregation, (2) Autonomous baseline budget execution, and (3) The cascading incentive distribution protocol.

\paragraph{Part 1: Initialization and Meso-Level State Aggregation}

Upon initialization, the simulation instantiates seven unique \texttt{BarangayAgent} objects—representing Liangan East, Esperanza, Poblacion, Binuni, Demologan, Mati, and Babalaya. Each object is dynamically loaded with distinct demographic constraints (e.g., population sizes ranging from 463 to 1,534 households) and localized Internal Revenue Allotment (IRA) budgets derived from municipal institutional audits \citep{Delena2025}.

Computationally, the Barangay Agent serves as the primary data aggregator for the Deep Reinforcement Learning (DRL) environment. During every simulation tick, the agent queries all constituent \texttt{HouseholdAgent} objects assigned to its jurisdiction. It calculates the localized compliance rate ($R_{bgy}$) in $O(N_b)$ time, where $N_b$ is the number of households in the specific barangay:
\begin{equation}
    R_{bgy} = \frac{\sum_{i=1}^{N_b} \mathbb{I}(\text{household}_i.\text{is\_compliant})}{N_b}
\end{equation}
This aggregated metric is crucial; it triggers the ``Habit Shield'' for the households (as detailed in Section 3.3.1) and is subsequently passed to the centralized \texttt{BacolodModel} to form the 17-dimensional State Space ($S_t$) evaluated by the AI Mayor \citep{TianReview2024}.

\paragraph{Part 2: Autonomous Base Layer Implementation (The Three Local Levers)}
While the Mayor Agent dynamically shifts municipal funds, the Barangay Agent continuously executes a static baseline policy using its localized budget allocation profile (e.g., 34\% IEC, 33\% Enforcement, 33\% Incentives). This ensures the community retains baseline governance even if the AI Mayor ignores them:

\begin{enumerate}
    \item \textbf{Base IEC (Micro-Interventions):} Barangay agents lack the financial capacity for mass media. Instead, their daily IEC budget is converted into grassroots activities using a greedy procurement heuristic. The algorithm prioritizes \textit{Recoridas} (roving announcements, $C_{rec} = \text{\textpeso} 200$) before allocating residual funds to \textit{Purok} meetings ($C_{pur} = \text{\textpeso} 1,000$). The resulting computational intensity applies a daily fractional scalar boost to the Attitude of all local households.
    \begin{equation}
        N_{\text{rec}} = \min\left(30, \left\lfloor \frac{Budget_{\text{IEC}}}{200} \right\rfloor \right)
        \label{eq:n_rec}
    \end{equation}

    \begin{equation}
        N_{\text{pur}} = \min\left(3, \left\lfloor \frac{\max(0, Budget_{\text{IEC}} - (N_{\text{rec}} \times 200))}{1000} \right\rfloor \right)
        \label{eq:n_pur}
    \end{equation}

    \begin{equation}
        I_{\text{IEC}} = \min\left(1.0, \ 0.5\left(\frac{N_{\text{rec}}}{30}\right) + 0.5\left(\frac{N_{\text{pur}}}{3}\right)\right)
        \label{eq:i_iec}
    \end{equation}

    \noindent \textit{where:}
    \begin{itemize}
        \item $N_{\text{rec}}$: The number of \textit{Recoridas} (roving grassroots announcements) funded and scheduled by the barangay for the given quarter.
        \item $N_{\text{pur}}$: The number of \textit{Purok} meetings (neighborhood assemblies) funded and scheduled by the barangay.
        \item $Budget_{\text{IEC}}$: The localized financial budget (in Philippine Pesos) strictly allocated to the barangay for Information, Education, and Communication campaigns.
        \item $I_{\text{IEC}}$: The Base IEC Intensity. This is a fractional scalar ranging from $0.0$ to $1.0$ that represents the combined computational strength of the barangay's local educational campaigns. It is used to apply a daily boost to the intrinsic environmental attitude of local households.
        \item $200$ ($C_{\text{rec}}$): The financial cost to conduct a single \textit{Recorida} (e.g., fuel for the barangay patrol vehicle).
        \item $30$: The maximum target limit for \textit{Recoridas} per quarter (representing roughly one announcement every 3 days).
        \item $1000$ ($C_{\text{pur}}$): The financial cost to conduct a single \textit{Purok} meeting (e.g., snacks, printed materials, and logistical setup).
        \item $3$: The maximum target limit for \textit{Purok} meetings per quarter (representing one major meeting per month).
        \item $0.5$: The equal mathematical weight ($50\%$) assigned to both \textit{Recoridas} and \textit{Purok} meetings when computing the final $I_{\text{IEC}}$ scalar, ensuring neither method completely dominates the intensity calculation.
    \end{itemize}

    \item \textbf{Base Enforcement (Tanod Spawning):} To execute local enforcement, the Barangay calculates its maximum affordable headcount based on its local enforcement fund. It dynamically instantiates persistent \texttt{EnforcementAgent} objects (flagged as \texttt{is\_municipal = False}). These Tanods have a strict apprehension capacity of 1 household per patrol and suffer from a 90\% stochastic distraction rate, reflecting their multi-duty nature. Their presence computes an ``Ambient Fear'' variable ($Base = 0.30$) to maintain a baseline perception of risk \citep{Zhao2022}. Cost of Enforcements will be deducted from the barangays fund, where $N_{\mathrm{Enforcers}}$ is the number of base enforcers (tanods) deployed, $W_{\mathrm{Daily}}$ is the minimum daily wage of a personnel in contract of service in Region X and $30$ as the days in contract of service.
    \begin{equation}
        C_{\mathrm{Enf}} = (N_{\mathrm{Enforcers}} \times W_{\mathrm{Daily}} \times 30)
    \end{equation}
    \label{eq:cost_of_enforcement}
    \vspace{-20pt} 
    \item \textbf{Base Incentives (Tier 1 Ledger):} The local incentive fund serves as the primary financial clearinghouse for positive reinforcement \citep{Chen2023}. The baseline computational cost of incentives ($C_{\mathrm{Inc}}$) is determined by the Equation \ref{eq:cost_of_incentives}, where  $N_{\mathrm{Compliant}}$ are the number of compliant \texttt{HouseholdAgents} and $V_{\mathrm{Reward}}$ is a \textpeso 500 subsidy drawn directly from this local pool.
    \begin{equation}
        C_{\mathrm{Inc}} = (V_{\mathrm{Reward}} \times N_{\mathrm{Compliant}})
    \end{equation}
    \label{eq:cost_of_incentives}
\end{enumerate}

\vspace{-45pt} % Adjust this number (e.g., -10pt, -20pt) to your liking
\paragraph{Part 3: The Cascading Incentive Distribution Protocol}

A major architectural feature of the \texttt{BarangayAgent} is its role as a financial clearinghouse. When a compliant \texttt{HouseholdAgent} attempts to claim a monetary reward, it invokes the barangay's \texttt{redeem\_incentive(amount)} method.

To accurately simulate the ``stacking effect'' of polycentric governance \citep{Ostrom2010}, this method utilizes a \textit{cascading resource allocation tree}:
\begin{itemize}
    \item \textbf{Tier 1 (Local Funds):} The algorithm first attempts to deduct the \textpeso 500 reward from the \texttt{local\_incentive\_fund}. 
    \item \textbf{Tier 2 (Municipal Augmentation):} If local funds are exhausted, the method queries the overarching \texttt{MayorAgent}. If the AI Mayor has injected targeted augmentation funds ($A_{\mathrm{Inc}, B_i}$) into this specific barangay, the reward is deducted from the municipal pool.
    \item \textbf{Tier 3 (Budget Exhaustion):} If both budget tiers are empty, the method returns \texttt{False}. The household's claim is rejected, triggering a computational "Perceived Unfairness" penalty that degrades the household's future intrinsic motivation \citep{Chen2023}.
\end{itemize}

This tiered financial logic natively simulates the ``victim of success'' dynamic \citep{WorldBank2022}. If local compliance surges rapidly, the fixed municipal allocation is diluted among a larger pool of claimants, rapidly exhausting the budget and forcing the DRL agent to constantly adjust its policy to maintain systemic motivation.

\vspace{0.5cm} 
\begin{algorithm}[H]
\caption{The Barangay's Administrative and Fiscal Loop}
\label{alg:barangay_execution_english}
\begin{algorithmic}[1]
\State \textbf{Starting Information Needed:} 
\State \parbox[t]{0.9\linewidth}{$\rightarrow$ The Barangay's assigned demographic profile, local baseline budget, and a registry of all resident households.\strut}
\Statex
\State \textbf{Phase 1: Daily Aggregation and Micro-Interventions}
\State \parbox[t]{0.9\linewidth}{$\rightarrow$ \textit{Data Collection:} Query all local households to calculate the official Barangay Segregation Rate. Broadcast this rate to the AI Mayor.\strut}
\State \parbox[t]{0.9\linewidth}{$\rightarrow$ \textit{Grassroots IEC:} Spend the daily local education budget. Buy cheap \textit{Recoridas} first; if money remains, buy \textit{Purok} meetings. Apply the resulting motivation boost to all local citizens.\strut}
\State \parbox[t]{0.9\linewidth}{$\rightarrow$ \textit{Personnel Management:} Count the active Barangay Tanods on the grid. If the local budget allows for more personnel, spawn new Tanods and dispatch them into the neighborhood.\strut}
\Statex
\State \textbf{Phase 2: Processing Citizen Rewards (The Cascading Ledger)}
\If{a citizen actively segregates and requests a municipal reward}
    \If{the Barangay has its own local incentive funds remaining}
        \State \parbox[t]{0.85\linewidth}{$\rightarrow$ \textit{Approval:} Pay the citizen using local funds. The citizen remains highly motivated.\strut}
    \ElsIf{the Mayor has provided supplemental municipal funds to this Barangay}
        \State \parbox[t]{0.85\linewidth}{$\rightarrow$ \textit{Approval:} Pay the citizen using the Mayor's augmentation budget.\strut}
    \Else
        \State \parbox[t]{0.85\linewidth}{$\rightarrow$ \textit{Rejection (Budget Exhausted):} Deny the reward. The citizen feels cheated, triggering a permanent drop in their future willingness to segregate.\strut}
    \EndIf
\EndIf
\end{algorithmic}
\end{algorithm}

\subsection{The Mayor Agent: Macro-Level Intervention and DRL Execution}

The \texttt{MayorAgent} serves as the centralized executive entity within the simulation. Unlike the \texttt{BarangayAgents}, which operate continuously on static, hardcoded baseline budgets, the Mayor Agent is fully dynamic and serves as the physical actuator for the Heuristic-Guided Deep Reinforcement Learning (HuDRL) policy. To systematically detail its computational role, this section is divided into four parts: (1) The temporal decision epoch, (2) The budget actuation pipelines, (3) The political survival constraint, and (4) The closed-loop communication topology.

\paragraph{Part 1: The Macro-Level Decision Epoch (The AI Bridge)}

Computationally, the \texttt{MayorAgent} acts as the structural bridge between the simulated Agent-Based Model (ABM) environment and the external PyTorch-based neural network. To simulate the administrative realities of municipal governance, the agent does not evaluate policy daily. Instead, its \texttt{step()} function enforces a strict 90-day (quarterly) temporal epoch \citep{Akbarpour2021}. 

At the start of each quarter ($t \pmod{90} == 0$), the Mayor Agent pauses local micro-executions to observe the global state vector (aggregated compliance, average attitudes, remaining annual budget, and political capital). It queries the HuDRL algorithm, receiving a continuous action vector representing a normalized budget allocation strategy. \textit{(Note: The specific neural network architecture, state-space constraints, and triage heuristics governing how the AI calculates this vector are extensively detailed in Section 3.4)}. Once received, the Mayor Agent immediately translates these abstract floating-point values into physical municipal interventions.

\paragraph{Part 2: Budget Translation and Macro-Intervention Levers}
The AI's action vector is partitioned and executed across the same three policy levers utilized by the Barangays, but at a massively scaled municipal capacity \citep{Ibanez2022}:

\begin{enumerate}
    \item \textbf{Macro-IEC (Mass Media Campaigns):} When the AI allocates funds to a barangay's IEC sector, the Mayor bypasses grassroots methods and invests in mass media. The allocation ($B_{lgu\_iec}$) funds broad-reach radio broadcasts ($C_{rad} = \text{\textpeso} 500$) and large-scale assembly events ($C_{evt} = \text{\textpeso} 5,000$). The intensity of municipal IEC ($I_{macro}$) is bounded by target saturation levels:
    \begin{align}
        n_{rad} &= \min \left( T_{rad}, \left\lfloor \frac{B_{lgu\_iec}}{C_{rad}} \right\rfloor \right) \\
        rem_{lgu} &= \max \left( 0, B_{lgu\_iec} - (n_{rad} \times C_{rad}) \right) \\
        n_{evt} &= \min \left( T_{evt}, \left\lfloor \frac{rem_{lgu}}{C_{evt}} \right\rfloor \right) \\
        I_{macro} &= \min \left( 1.0, \omega_{rad} \left( \frac{n_{rad}}{T_{rad}} \right) + \omega_{evt} \left( \frac{n_{evt}}{T_{evt}} \right) \right)
    \end{align}

    \noindent \textit{Where:}
    \begin{itemize}
        \item $B_{lgu\_iec}$: The municipal budget allocation (in Philippine Pesos) injected by the AI Mayor specifically into a target barangay's IEC sector.
        \item $n_{rad}$: The calculated number of broad-reach radio broadcasts successfully funded by the municipal allocation.
        \item $rem_{lgu}$: The residual municipal budget remaining after funding the target number of radio broadcasts, which cascades down to fund larger events.
        \item $n_{evt}$: The calculated number of large-scale assembly events successfully funded using the residual municipal budget.
        \item $I_{macro}$: The Macro-Level IEC Intensity. A fractional scalar (ranging from $0.0$ to $1.0$) representing the combined computational impact and reach of the municipal mass media campaign.
        \item $A$ and $SN$: The intrinsic Attitude and Subjective Norms of the targeted households, which receive an immediate artificial boost from the $I_{macro}$ intervention.
        \item $500$ ($C_{rad}$): The financial cost (in Philippine Pesos) to air a single broad-reach radio broadcast.
        \item $5000$ ($C_{evt}$): The financial cost (in Philippine Pesos) to organize and host a single large-scale municipal assembly event.
        \item $T_{rad}$: The target saturation limit for radio broadcasts, representing the maximum effective number of broadcasts per quarter before encountering diminishing returns.
        \item $T_{evt}$: The target saturation limit for large-scale assembly events per quarter.
        \item $\omega_{rad}$ and $\omega_{evt}$: The mathematical weights assigned to the radio broadcasts and assembly events, respectively. These dictate their proportional influence when calculating the final overall intensity ($I_{macro}$).
    \end{itemize}

    This intensity instantly artificially boosts the Attitude ($A$) and Social Norms ($SN$) of households in the targeted jurisdiction.
    
    \item \textbf{Macro-Enforcement (LGU Strike Force) and Object Destruction:} To execute localized crackdowns, enforcement funds are translated into elite ``Strike Force'' personnel. Unlike local Tanods, these agents are hired as temporary shock troops. Incorporating the Region X Daily Minimum Wage of \textpeso 400 dictated that deploying a single municipal enforcer on a strict 30-day temporal contract consumed exactly \textpeso 12,000. The municipal enforcement cost ($C_{\mathrm{Enf}}$) is linearly tied to human capital:
    \begin{equation}
        C_{\mathrm{Enf}} = (N_{\mathrm{Enforcers}} \times W_{\mathrm{Daily}} \times 30)
    \end{equation}
    The Mayor dynamically instantiates these high-capacity \texttt{EnforcementAgents} (flagged as \texttt{is\_municipal = True}) and injects them into the grid. They bypass the distraction filter (achieving a 100\% daily patrol rate), execute up to 25 apprehensions per tick, and issue severe \textpeso 1,000 citations \citep{Wang2023}. 
    
    Crucially, these agents are bound by their strict 30-day limit enforced via computational object destruction. During each daily tick, the agent decrements its internal \texttt{contract\_days} integer. Upon reaching zero, the object invokes a self-destruct sequence (\texttt{self.model.grid.remove\_agent()}), physically deleting itself from the spatial grid. This mechanical constraint prevents the AI Mayor from permanently locking down a barangay with a single payment. While the Strike Force generates an immediate spike in compliance via the "Fear Override," their sudden deletion guarantees a rapid behavioral collapse in subsequent ticks unless the AI learns to transition the community toward intrinsic motivation (IEC) or expends further budget in the next quarter.
    
    \item \textbf{Macro-Incentives (Economic Augmentation):} Funds allocated to positive reinforcement are channeled directly into the target \texttt{BarangayAgent}'s Tier 2 financial ledger. This ensures the barangay has the liquidity to process citizen reward claims without exhausting their own constrained local funds.
\end{enumerate}

\paragraph{Part 3: Political Capital Maintenance (Survival Mechanics)}

To prevent the AI from adopting a computationally optimal but socially disastrous draconian strategy, the Mayor Agent is bound by a "Political Capital" constraint \citep{Hertweck2023}. 

In the model, every municipal fine (exceeding \textpeso 500) levied by the Strike Force triggers an event-driven interrupt that applies a fractional deduction to the Mayor's global public approval rating ($\Delta PC = -0.00010$). While municipal enforcement guarantees an immediate spike in segregation via the "Fear Override", it rapidly burns through this capital. If the Mayor's Political Capital drops below the critical threshold of $0.10$, the \texttt{BacolodModel} forces a premature termination of the simulation instance. This simulates a "Political Collapse" due to extreme public backlash, aggressively penalizing the AI during the training phase.

\paragraph{Part 4: Inter-Agent Communication Topology}

The resulting multi-level architecture forms a closed, bi-directional computational feedback loop, officially grounding the ABM as a Markov Decision Process (MDP) \citep{Kompella2020}:
\begin{itemize}
    \item \textbf{Bottom-Up Aggregation:} At the micro-level, Household Agents calculate their utility and update their binary compliance state. The Barangay Agents act as meso-level data nodes, aggregating these individual states into localized metrics ($R_{comp}$) and passing them to the centralized environment.
    \item \textbf{Top-Down Actuation:} The Mayor Agent observes this environment, consults the external DRL policy, and deploys financial capital downward. This capital physically alters the simulation grid (spawning new enforcement objects) or directly mutates the internal psychological variables of the households, forcing a new behavioral reaction for the next decision epoch.
\end{itemize}

\vspace{0.5cm} 
\begin{algorithm}[H]
\caption{The Mayor Agent's Quarterly Intervention Cycle}
\label{alg:mayor_decision_english}
\begin{algorithmic}[1]
\State \textbf{Starting Information Needed:} 
\State \parbox[t]{0.9\linewidth}{$\rightarrow$ The current global segregation rate, remaining annual budget, and political approval rating.\strut}
\Statex
\State \textbf{Step 1: The AI Consultation (Occurs every 90 Days)}
\State \parbox[t]{0.9\linewidth}{$\rightarrow$ Pause the simulation and transmit the city's current status to the HuDRL Artificial Intelligence.\strut}
\State \parbox[t]{0.9\linewidth}{$\rightarrow$ Receive a finalized, mathematically optimized budget blueprint allocating funds across the 7 barangays.\strut}
\Statex
\State \textbf{Step 2: Execute Budget Actuations (Top-Down Flow)}
\For{every Barangay in the municipality}
    \State \parbox[t]{0.88\linewidth}{$\rightarrow$ \textit{Information Campaigns:} If funds are allocated to IEC, purchase radio airtime and fund large events. Instantly broadcast a motivation boost to all households in this barangay.\strut}
    \State \parbox[t]{0.88\linewidth}{$\rightarrow$ \textit{Strike Force Deployment:} If funds are allocated to Enforcement, calculate how many elite inspectors can be hired. Spawn them into the barangay's grid on a strict 30-day contract.\strut}
    \State \parbox[t]{0.88\linewidth}{$\rightarrow$ \textit{Incentive Augmentation:} If funds are allocated to Rewards, deposit the money directly into the local Barangay's ledger to subsidize citizen payouts.\strut}
\EndFor
\Statex
\State \textbf{Step 3: Monitor Survival Constraints}
\If{the Mayor's Political Capital (Approval Rating) $\le 10\%$}
    \State \parbox[t]{0.88\linewidth}{$\rightarrow$ Trigger an immediate simulation halt. The administration has collapsed due to excessive punitive measures and public resentment.\strut}
\Else
    \State \parbox[t]{0.88\linewidth}{$\rightarrow$ Resume the simulation. Allow households and enforcers to interact with the newly altered environment for the next 90 days.\strut}
\EndIf
\end{algorithmic}
\end{algorithm}

\subsection{Formalization as a Markov Decision Process}

To facilitate the autonomous discovery of optimal strategies in the absence of longitudinal historical data, the interaction between the Mayor Agent (representing the municipal Local Government Unit) and the household agents was formalized as a finite-horizon Markov Decision Process (MDP). This mathematical framework transitioned the Agent-Based Model from a descriptive simulation into a prescriptive optimization environment, where the Mayor Agent iteratively interacted with a stochastic state space to maximize a cumulative utility function \citep{Kompella2020}. The MDP is formally defined by the quintuple $\mathcal{M} = \langle S, A, P, R, \gamma \rangle$, representing the state space, action space, transition probability, reward signal, and discount factor, respectively.

The State Space ($S$) is characterized as a multi-dimensional, continuous vector $S_t \in \mathbb{R}^{17}$, which provides the Mayor Agent with a comprehensive snapshot of the municipal environment at each decision epoch. This vector encompasses the mean compliance rates across the seven constituent barangays, the normalized remaining municipal budget for the fiscal year, and the simulation time index to account for seasonal or temporal variations in waste generation. By including these macroscopic variables, the Mayor Agent can perceive the heterogeneous social landscape of the municipality, allowing for localized rather than monolithic policy responses \citep{HaMinh2025}. 

Correspondingly, the Action Space ($A$) is defined as a continuous vector $A_t \in \mathbb{R}^{21}$ belonging exclusively to the Mayor Agent. This high-dimensional space represents the granular allocation of the quarterly municipal intervention budget across the three primary policy levers—Enforcement, Incentives, and Information Education Campaigns (IEC)—for each individual barangay. Crucially, the \textit{BarangayAgents'} base budgets are executed deterministically; it is the Mayor Agent that dynamically controls the supplemental funds. The choice of a continuous action space over a discrete one is critical for achieving "policy precision," as it enables the Mayor Agent to fine-tune resource intensity to the specific tipping points of household behavioral change \citep{Rajesh2025}.

The system dynamics are governed by the Transition Function ($P$), denoted as $P(S_{t+1} | S_t, A_t)$. Within this framework, the ABM serves as the implicit transition engine; unlike classical control problems where transitions are defined by explicit differential equations, the state evolution here emerges from the aggregate, non-linear behaviors of heterogeneous household agents updating their internal utility based on the Theory of Planned Behavior (TPB). Consequently, the simulation acts as a generative model where the subsequent state $S_{t+1}$ is a stochastic manifestation of the community’s collective response to both the barangay baseline and the Mayor Agent's active policy intervention $A_t$ \citep{TianReview2024, Jimenez2025}. To evaluate the efficacy of these interventions, the Reward Function ($R$) provides a scalar feedback signal $R_t$ that balances the Mayor Agent’s multi-objective mandates. This function is designed to reward high aggregate compliance while penalizing fiscal overspending and the social friction associated with excessive punitive enforcement \citep{Jimenez2025}. Finally, a Discount Factor of $\gamma = 0.99$ is applied to the objective function, ensuring that the Mayor Agent prioritizes strategies that foster long-term, sustained behavioral shifts rather than pursuing ephemeral, short-term compliance spikes that may deplete the annual municipal budget prematurely.